\section{Overview}
\label{sec:overview}

A user writes a grammar as an ordinary C++ type.  The compiler lowers it,
entirely at compile time, to a data-parallel scanner whose structure is fixed by
the scan algebra $\SA$ (Figure~\ref{fig:framework}).  This section gives the intuition; the
formalization follows in \S\ref{sec:model}.

% Framework pipeline diagram (placeholder; may be hand-drawn).
\begin{figure}[t]
\centering
\resizebox{\linewidth}{!}{%
\begin{tikzpicture}[
  node distance=6mm,
  box/.style={draw, rounded corners, align=center, inner sep=3pt, font=\small},
  prim/.style={draw, rounded corners, fill=black!4, align=center, inner sep=2pt, font=\scriptsize},
  ar/.style={-{Latex[length=1.6mm]}, thick}]

\node[box] (g) {\textbf{Grammar}\\[-1pt]\scriptsize C++ types};
\node[box, right=of g] (a) {\textbf{Analysis}\\[-1pt]\scriptsize refs, recursion};
\node[box, right=of a] (c) {\textbf{Classifier}\\[-1pt]\scriptsize regular/mode/dyck/table};
\node[box, right=of c] (s) {\textbf{Schedule}\\[-1pt]\scriptsize descriptors $\to$ stages};
\node[box, right=of s] (b) {\textbf{Backends}\\[-1pt]\scriptsize CPU / CUDA};

\draw[ar] (g) -- (a);
\draw[ar] (a) -- (c);
\draw[ar] (c) -- (s);
\draw[ar] (s) -- (b);

\node[prim, below=8mm of s, xshift=-18mm] (r1) {R1 classify};
\node[prim, right=of r1] (r4) {R4 monoid scan};
\node[prim, right=of r4] (r5) {R5 compact};
\node[prim, below=2mm of r1] (r3) {R3 select};
\node[prim, right=of r3] (r6) {R6 nest reduce};
\node[prim, right=of r6] (r7) {R7 assoc reduce};
\draw[ar, dashed] (s) -- (r4);
\end{tikzpicture}%
}
\caption{The compiler and scan algebra $\SA$.  A grammar is lowered, entirely at
compile time, to a schedule of primitives executed by CPU/CUDA backends.}
\label{fig:framework}
\end{figure}


\paragraph{The surface language.}
The whole user input is one C++ type: rules are tag structs, rule bodies are
type aliases, and there is no run-time grammar object.  Figure~\ref{fig:grammar}
is a JSON fragment.

\begin{figure}[t]
\begin{lstlisting}
struct json_grammar {
  struct value {}; struct string {};
  struct number {}; struct ws {};
  using rules         = type_list<value, string,
                                  number, ws>;
  using lexical_rules = type_list<string, number, ws>;
  using start = value;
  template <class> struct def;
};
template <> struct json_grammar::def<
    json_grammar::string> {
  using type = seq<lit<'"'>,
      star<alt<seq<lit<'\\'>, any>,
               seq<notp<cls<'"','\\'>>, any>>>,
      lit<'"'>>;
};
\end{lstlisting}
\caption{A grammar is a type.  \code{lexical\_rules} marks the rules that mask;
all other rules are classified as structural or recursive by the front end
(automaton + FIRST/LAST + the bracket guard), never by pattern matching.
The full seven-grammar suite ships in the artifact.}
\label{fig:grammar}
\end{figure}

\begin{table}[t]
\centering\small
\caption{Surface constructs and their grammar meaning.}
\label{tab:dsl}
\begin{tabular}{@{}ll@{}}
\toprule
Construct & Grammar meaning \\
\midrule
\code{lit<c>}, \code{cls<c...>}, \code{any} & terminal, byte class, wildcard \\
\code{seq<>}, \code{alt<>} & concatenation, alternation \\
\code{opt<>}, \code{star<>}, \code{plus<>} & Kleene operators \\
\code{rule<T>} & nonterminal reference \\
\code{notp<T>} & negated predicate (lookahead) \\
\code{rules}, \code{start}, \code{def<T>} & rule set, start symbol, body \\
\code{lexical\_rules} & mask these rules \\
\code{fallback\_rules} & opt-out to a fallback backend \\
\bottomrule
\end{tabular}
\end{table}

\paragraph{From a branch program to masked data flow.}
Consider the scalar loop that masks JSON strings:
\begin{lstlisting}
in = esc = false;
for (i = 0; i < n; ++i) {
  c = data[i];
  if (in) {
    if (esc)              esc = false;
    else if (c == '\\')   esc = true;
    else if (c == '"')    in  = false;
  } else {
    if (c == '"')         in  = true;
    else if (structural(c)) emit(i);
  }
}
\end{lstlisting}
Every condition is a predicate on the current byte and a bounded register.  The
rewrite replaces each byte by the state transform it induces and the loop by a
\emph{prefix scan} over the transformation monoid; a branch becomes a masked
select.  For the 1-bit \texttt{in}-flag the transform is affine,
$x'=(a\wedge x)\oplus b$, and its prefix is a single
\textsf{prefix\_xor} over the quote mask once escaped quotes are removed.

\paragraph{Masking from an automaton, not a shape.}
We do not pattern-match rule syntax.  Each lexical rule is compiled to a DFA by
Brzozowski derivatives, and a byte is masked iff the DFA transition it takes is
one that \emph{consumes it as part of an in-progress token}.  This one rule
subsumes prefix-escaped strings, doubled-quote strings, line and block comments,
and multi-mode languages.  Specializations (affine, doubled, comment) are
emitted only when an exhaustive check shows them equivalent to the automaton.

\paragraph{Parsing as well-nested reduction.}
For a well-nested (visibly pushdown) recursion, the stack is replaced by a
depth group-scan plus a sort: depth is the prefix sum over open/close brackets,
and matched pairs are adjacent tokens after a stable sort by depth.  This is
\textsf{cuJSON}'s stack-free matching, derived from the grammar rather than
special-cased.

\paragraph{Boundary, kernels, chunking --- not transfer.}
The framework owns the device boundary: \code{scan\_device} is zero-copy
(input on device, results on device), the \code{scan\_host} adapter is the only
place with bulk transfers, and \code{scan\_chunk} carries the prefix-monoid
state and output offset so chunks can be pipelined.  This separation is what
makes the three throughput tiers of \S\ref{sec:eval} meaningful: a compute
roofline, an on-device number, and a host-boundary number.
